% 第6章 低维中的失序
\chapter{低维中的失序}

\section{自发破缺对称}

自发破缺对称是统计力学与粒子物理中的普遍现象。它可以在有限温度下发生，此时序参量与哈密顿量不对易。为使讨论集中，考虑 $z$ 方向的自旋密度波，其算符定义为

\begin{equation*}
S _ {\mathbf {q}} ^ {z} = \sum_ {i} e ^ {i \mathbf {q} \cdot \mathbf {x} _ {i}} S _ {i} ^ {z} ,
\tag{6.1}
\end{equation*}

$i$ 标记格点，$\mathbf{q}$ 是序波矢。

用序场 $h$ 给转动不变的哈密顿量 $H^{0}$ 加上对称破缺项：

\begin{equation*}
\mathcal {H} ( h ) = \mathcal {H} _ {0} - h S _ {\mathbf {q}} ^ {z}.
\tag{6.2}
\end{equation*}

每格点磁化为

\begin{align*}
m _ {\mathbf {q}} ( h ) & = \frac {1}{\mathcal {N} Z} \mathrm{Tr} \left[ e ^ {- \mathcal {H} ( h ) / T} S _ {\mathbf {q}} ^ {z} \right] ,\\
Z & = \mathrm{Tr} \left[ e ^ {- \mathcal {H} ( h ) / T} \right].
\tag{6.3}
\end{align*}

为方便起见，把讨论限于关于原点反射对称的格子，此时 (6.3) 中的 $S_{q}^{z}$ 厄米，$m_{q}$ 为实。

有限场下 $m_{\mathbf{q}}(h) > 0$，因为磁化由序场诱导。

\medskip
\noindent\textbf{定义 6.1}\quad 若在热力学极限下，即使序场从正侧趋于零，体系仍维持有限磁化，则称体系具有自发破缺对称：

\begin{equation*}
\lim _ {h \to 0 ^ {+}} \lim _ {\mathcal {N} \to \infty} m _ {\mathbf {q}} ( h, \mathcal {N}, T ) \neq 0.
\tag{6.4}
\end{equation*}
\medskip

(6.4) 中极限的次序至关重要。通过考察两点关联函数，可以在没有序场的情况下发现自发破缺对称：

\begin{equation*}
S ^ {\alpha \alpha} ( \mathbf {q} ) = \lim _ {h \to 0 ^ {+}} \frac {1}{Z \mathcal {N}} \mathrm{Tr} \left[ e ^ {- \mathcal {H} ( h ) / T} S _ {\mathbf {q}} ^ {\alpha} S _ {- \mathbf {q}} ^ {\alpha} \right] , \quad \alpha \in \{ x, y, z \}.
\tag{6.5}
\end{equation*}

没有序场时，关联在自旋空间中各向同性，即与方向 $\alpha$ 无关。可以证明（见习题）：动量 $\mathbf{q}$ 处的自发破缺对称蕴含两点关联函数中的真长程序：

\begin{equation*}
\lim _ {\mathcal {N} \rightarrow \infty} \left[ \mathcal {N} ^ {- 1} S ^ {\alpha \alpha} ( \mathbf {q} ) \right] > 0 ,
\tag{6.6}
\end{equation*}

或者等价地

\begin{equation*}
\lim _ {| \mathbf {x} _ {i} - \mathbf {x} _ {j} | \to \infty} \; \lim _ {\mathcal {N} \to \infty} \langle \mathbf {S} _ {i} \cdot \mathbf {S} _ {j} \rangle \neq 0 ,
\tag{6.7}
\end{equation*}

其中 $|x_{i}-x_{j}|$ 是格点 $i$、$j$ 之间的距离。“真长程序”一词用于把 (6.7) 与“准长程序”区分开，后者指关联在大距离上按幂律衰减。

\section{Mermin--Wagner 定理}

预期对具有有序基态的体系，热激发会在有限温度下削弱自旋关联；经典与量子体系皆然。当温度远高于典型耦合能量尺度 $J$ 时，预期自旋在远距离不关联，且没有序场时 $m_{q}$ 为零。这要求在有序相与无序相之间存在某个相变温度 $T_{c}$。

Heisenberg 模型的有序基态破缺自旋空间的连续 $O(3)$ 对称。把序场的取向选在球面任意方向，就能让自发磁化指向那里。连续对称的后果在低维格子上特别显著。事实上我们将会看到，在一维和二维，相变恰好位于 $T=0$：无穷小的温度就会使热激发破坏自旋的有序。

1966 年，Hohenberg 利用 Bogoliubov 不等式（下文定义）证明了一维、二维中有限温超流的不存在（见文献注释）。Mermin 与 Wagner 沿用实质相同的方法，证明一维与二维的短程自旋模型在任何有限温度下都不能自发有序。下面给出的证明专门针对量子 Heisenberg 模型，但容易推广到更大一类具有连续对称性的模型。

\medskip
\noindent\textbf{Mermin--Wagner 定理 6.2}\quad 对量子 Heisenberg 模型

\begin{equation*}
\mathcal {H} = \frac {1}{2} \sum_ {ij} J _ {ij} \, \mathbf {S} _ {i} \cdot \mathbf {S} _ {j} - h S _ {\mathbf {q}} ^ {z}
\tag{6.8}
\end{equation*}

若其短程相互作用满足

\begin{equation*}
\bar {J} = \frac {1}{2 \mathcal {N}} \sum_ {ij} \left| J _ {ij} \right| \left| \mathbf {x} _ {i} - \mathbf {x} _ {j} \right| ^ {2} < \infty ,
\tag{6.9}
\end{equation*}

则在一维和二维，有限温度下不存在自发破缺的自旋对称：

\begin{equation*}
\lim _ {h \to 0 ^ {+}} \lim _ {\mathcal {N} \to \infty} m _ {\mathbf {q}} ( h, \mathcal {N} ) = 0.
\tag{6.10}
\end{equation*}
\medskip

任意两个算符 $A$、$B$ 之间的标量积定义为

\begin{equation*}
( A, B ) = \frac {1}{Z} \sum_{n, m} ' \langle n | A ^ {\dagger} | m \rangle \langle m | B | n \rangle \left( \frac {e ^ {- E _ {m} / T} - e ^ {- E _ {n} / T}}{E _ {n} - E _ {m}} \right),
\tag{6.11}
\end{equation*}

$\{E_{n},|n\rangle\}$ 是 $\mathcal{H}$ 的能量与本征态；$\Sigma'$ 排除 $E_{n}=E_{m}$ 的项。容易看出括号中的权重非负，因此 (6.11) 是乘积 Hilbert 空间 $\{|n\rangle\}\times\{|n\rangle\}$ 中的正规标量积。$A=B$ 时，(6.11) 是 $A$ 的平方范数，且由 (B.27) 与 (B.10)，它是算符 $A$ 的磁化率：

\begin{equation*}
( A, A ) = \operatorname{Re} R _ {AA} ( 0 ) = \chi_ {AA}.
\tag{6.12}
\end{equation*}

用不等式 $\tanh(x) < x$ 容易验证

\begin{equation*}
0 < \left( \frac {e ^ {- E _ {m} / T} - e ^ {- E _ {n} / T}}{E _ {n} - E _ {m}} \right) \leq \frac {1}{2 T} \left( e ^ {- E _ {m} / T} + e ^ {- E _ {n} / T} \right).
\tag{6.13}
\end{equation*}

于是得到不等式

\begin{equation*}
( A, A ) \leq \frac {1}{2 T} \langle A ^ {\dagger} A + A A ^ {\dagger} \rangle .
\tag{6.14}
\end{equation*}

对标量积用 Cauchy--Schwarz 不等式并结合 (6.14)：

\begin{align*}
\left| ( A, B ) \right| ^ {2} & \leq ( A, A ) ( B, B )\\
& \leq \frac {1}{2 T} \langle A ^ {\dagger} A + A A ^ {\dagger} \rangle \, ( B, B ).
\tag{6.15}
\end{align*}

定义算符 $C$ 使

\begin{equation*}
B \equiv [ C ^ {\dagger}, \mathcal {H} ].
\tag{6.16}
\end{equation*}

则

\begin{align*}
( A, B ) & = \langle [ C ^ {\dagger}, A ^ {\dagger} ] \rangle ,\\
( B, B ) & = \langle [ C ^ {\dagger}, [ \mathcal {H}, C ] ] \rangle ,
\tag{6.17}
\end{align*}

由 (6.15) 即得 Bogoliubov 不等式：

\begin{equation*}
\left| \langle [ C ^ {\dagger}, A ^ {\dagger} ] \rangle \right| ^ {2} \leq \frac {1}{2 T} \langle A ^ {\dagger} A + A A ^ {\dagger} \rangle \, \langle [ C ^ {\dagger}, [ \mathcal {H}, C ] ] \rangle .
\tag{6.18}
\end{equation*}

现在取算符

\begin{equation*}
C = S _ {\mathbf {k}} ^ {x} \, , \quad A = S _ {- \mathbf {k} - \mathbf {q}} ^ {y} ,
\tag{6.19}
\end{equation*}

这给出

\begin{align*}
\frac {1}{2} \langle A ^ {\dagger} A + A A ^ {\dagger} \rangle & = \langle S _ {\mathbf {k} + \mathbf {q}} ^ {y} S _ {- \mathbf {k} - \mathbf {q}} ^ {y} \rangle\\
& = \mathcal {N} S ^ {yy} ( \mathbf {k} + \mathbf {q} ) .
\tag{6.20}
\end{align*}

$S^{yy}$ 是 (6.5) 定义的等时关联函数。

把 (6.19) 用于 (6.18) 左端：

\begin{align*}
\langle [ C ^ {\dagger}, A ^ {\dagger} ] \rangle & = \langle [ S _ {- \mathbf {k}} ^ {x}, S _ {\mathbf {k} + \mathbf {q}} ^ {y} ] \rangle\\
& = i \langle S _ {\mathbf {q}} ^ {z} \rangle = i \mathcal {N} m _ {\mathbf {q}} .
\tag{6.21}
\end{align*}

$F(\mathbf{k})$ 是双对易子函数

\begin{equation*}
F ( \mathbf {k} ) = \mathcal {N} ^ {- 1} \left\langle \left[ S _ {- \mathbf {k}} ^ {x}, [ \mathcal {H}, S _ {\mathbf {k}} ^ {x} ] \right] \right\rangle .
\tag{6.22}
\end{equation*}

于是 Bogoliubov 不等式 (6.18) 化为

\begin{equation*}
m _ {\mathbf {q}} ^ {2} \leq \frac {1}{T} S ^ {yy} ( \mathbf {k} + \mathbf {q} ) \, F ( \mathbf {k} ).
\tag{6.23}
\end{equation*}

用哈密顿量 (6.8) 可以对 $F(\mathbf{k})$ 作如下估计：

\begin{align*}
F ( \mathbf {k} ) & = h m _ {\mathbf {q}} + i \mathcal {N} ^ {- 1} \sum_ {ijl} e ^ {i \mathbf {k} ( \mathbf {x} _ {l} - \mathbf {x} _ {i} )} J _ {jl} \left\langle \left[ S _ {i} ^ {x}, \left( S _ {j} ^ {z} S _ {l} ^ {y} - S _ {j} ^ {y} S _ {l} ^ {z} \right) \right] \right\rangle\\
& = h m _ {\mathbf {q}} + \mathcal {N} ^ {- 1} \sum_ {jl} J _ {jl} \left( \cos [ \mathbf {k} ( \mathbf {x} _ {j} - \mathbf {x} _ {l} ) ] - 1 \right) \left\langle S _ {j} ^ {y} S _ {l} ^ {y} + S _ {j} ^ {z} S _ {l} ^ {z} \right\rangle\\
& \leq h m _ {\mathbf {q}} + \frac {| \mathbf {k} | ^ {2}}{2 \mathcal {N}} \sum_ {jl} \left| J _ {jl} \right| \left| \mathbf {x} _ {j} - \mathbf {x} _ {l} \right| ^ {2} \left| \langle S _ {j} ^ {y} S _ {l} ^ {y} + S _ {j} ^ {z} S _ {l} ^ {z} \rangle \right| .
\tag{6.24}
\end{align*}

考察 $(\mathbf{S}_{j} + \mathbf{S}_{l})^{2}$ 与 $S_{i}^{2}$ 的最大本征值，得上界

\begin{equation*}
\left| \langle S _ {j} ^ {y} S _ {l} ^ {y} + S _ {j} ^ {z} S _ {l} ^ {z} \rangle \right| \leq \left| \langle \mathbf {S} _ {j} \cdot \mathbf {S} _ {l} \rangle \right| \leq S ( S + 1 )
\tag{6.25}
\end{equation*}

于是

\begin{equation*}
F ( \mathbf {k} ) \leq h m _ {\mathbf {q}} + S ( S + 1 ) \bar {J} \, | \mathbf {k} | ^ {2} ,
\tag{6.26}
\end{equation*}

其中 $\bar{J}$ 已在 (6.9) 定义。

双对易子函数 $F(\mathbf{k})$ 含有动力学信息，因为它显式依赖 $\mathcal{H}$。由 (6.26)，只有当 $\bar{J}$ 有限——即相互作用在大距离上衰减得足够快——时，$F(\mathbf{k})$ 才有有限上界。

组合 (6.20)、(6.21)、(6.23) 与 (6.26)：

\begin{equation*}
m _ {\mathbf {q}} ^ {2} \leq \frac {1}{T} S ^ {yy} ( \mathbf {k} + \mathbf {q} ) \left[ h m _ {\mathbf {q}} + S ( S + 1 ) \bar {J} | \mathbf {k} | ^ {2} \right]
\tag{6.27}
\end{equation*}

或

\begin{equation*}
S ^ {yy} ( \mathbf {k} + \mathbf {q} ) \geq \frac {T m _ {\mathbf {q}} ^ {2}}{h m _ {\mathbf {q}} + S ( S + 1 ) \bar {J} | \mathbf {k} | ^ {2}}.
\tag{6.28}
\end{equation*}

$S^{yy}(\mathbf{k})$ 的和受自旋大小限制：

\begin{equation*}
\mathcal {N} ^ {- 1} \sum_ {\mathbf {k}} S ^ {yy} ( \mathbf {k} + \mathbf {q} ) = \mathcal {N} ^ {- 1} \sum_ {i} \langle \left( S _ {i} ^ {y} \right) ^ {2} \rangle \leq S ( S + 1 ).
\tag{6.29}
\end{equation*}

用 (6.29) 对 (6.28) 两端求 $\mathbf{k}$ 和：

\begin{equation*}
S ( S + 1 ) \geq \frac {T m _ {\mathbf {q}} ^ {2}}{( 2 \pi ) ^ {d}} \int_ {0} ^ {\bar {k}} \frac {d k \, k ^ {d - 1}}{h m _ {\mathbf {q}} + S ( S + 1 ) \bar {J} | \mathbf {k} | ^ {2}} ,
\tag{6.30}
\end{equation*}

这里已把动量和换成积分，并取 $\bar{k}$ 小于布里渊区边界波矢以得到右端的解析估计。一维情形得

\begin{align*}
S ( S + 1 ) & \geq \frac {T m _ {\mathbf {q}} ^ {2}}{2 \pi \sqrt {\bar {J} S ( S + 1 ) h m _ {\mathbf {q}}}} \arctan \left[ \bar {k} \sqrt { \bar {J} S ( S + 1 ) / h m _ {\mathbf {q}} } \right]\\
& \xrightarrow[h \rightarrow 0]{} c \, \frac {T m _ {\mathbf {q}} ^ {3/2}}{\sqrt {h \bar {J} S ( S + 1 )}} \quad ( d = 1 ) ,
\tag{6.31}
\end{align*}

$c$ 为数值常数。反解 (6.31)：

\begin{equation*}
m _ {\mathbf {q}} \leq c \, \frac {S ( S + 1 ) \bar {J} ^ {1/3}}{T ^ {2/3}} \, h ^ {1/3}.
\tag{6.32}
\end{equation*}

二维情形的类似分析给出

\begin{equation*}
S ( S + 1 ) \geq \frac {T m _ {\mathbf {q}} ^ {2}}{4 \pi S ( S + 1 ) \bar {J}} \ln \left( 1 + \frac {\bar {J} S ( S + 1 ) \bar {k} ^ {2}}{h m _ {\mathbf {q}}} \right) ,
\tag{6.33}
\end{equation*}

反解得界

\begin{equation*}
m _ {\mathbf {q}} \leq c \, \frac {S ( S + 1 ) \bar {J} ^ {1/2}}{T ^ {1/2}} \, \frac {1}{\sqrt {| \ln h |}}.
\tag{6.34}
\end{equation*}

方程 (6.32) 与 (6.34) 证明：对一切有限的 $T$、$S$、$\bar{J}$，磁化必须随序场趋于零。由于这些界不依赖格子大小，结果在热力学极限下成立。证毕。

由于定理 6.2 对所有 $S$ 成立，可以推断它对 Heisenberg 模型的经典（与 $S$ 无关）极限同样成立：

\begin{equation*}
\mathcal {H} ^ {cl} \equiv \frac {1}{2} \sum_ {ij} J _ {ij} ^ {cl} \, \hat {\Omega} _ {i} \cdot \hat {\Omega} _ {j} - h ^ {cl} \hat {\Omega} _ {\mathbf {q}} ^ {z} ,
\tag{6.35}
\end{equation*}

其中 $\hat{\Omega}$ 是 $c$ 数单位矢量。(6.8) 与 (6.35) 的对应由参数标度给出：

\begin{align*}
J _ {ij} S ( S + 1 ) & \to J ^ {cl}\\
h S & \to h ^ {cl}\\
m _ {\mathbf {q}} / S & \to m _ {\mathbf {q}} ^ {cl}.
\tag{6.36}
\end{align*}

用这些参数即可从不等式 (6.32) 与 (6.34) 中消去 $S$。因此定理 6.2 直接适用于经典模型 (6.35)。证毕。

在 $d=3$，(6.30) 中的积分在 $h \rightarrow 0$ 时并不发散。因此可以在某个有限温度下满足不等式。于是预期三维中低温下存在自发对称破缺与有序相。

\section{$T = 0$ 的量子失序}

Mermin--Wagner 定理在 $T = 0$ 不适用。设在无穷小序场下存在唯一基态 $|0\rangle$。$T = 0$ 时任何算符的期望值为

\begin{equation*}
\langle A \rangle = \langle 0 | A | 0 \rangle .
\tag{6.37}
\end{equation*}

标量积 (6.11) 成为

\begin{equation*}
( A, B ) = \operatorname{Re} \sum_ {m \neq 0} \frac {\langle 0 | A ^ {\dagger} | m \rangle \langle m | B | 0 \rangle}{E _ {m} - E _ {0}}.
\tag{6.38}
\end{equation*}

按 (6.15) 与 (6.19) 取算符 $A, B, C$。磁化率为

\begin{align*}
\chi^ {yy} ( \mathbf {k} ) & = \mathcal {N} ^ {- 1} \left( S _ {- \mathbf {k}} ^ {y}, S _ {- \mathbf {k}} ^ {y} \right)\\
& = \mathcal {N} ^ {- 1} \sum_ {m \neq 0} \frac {\left| \langle 0 | S _ {\mathbf {k}} ^ {y} | m \rangle \right| ^ {2} + \left| \langle 0 | S _ {- \mathbf {k}} ^ {y} | m \rangle \right| ^ {2}}{E _ {m} - E _ {0}} .
\tag{6.39}
\end{align*}

把 (6.39) 代入 (6.15) 第一行并结合 (6.26)，得零温 Bogoliubov 不等式：

\begin{align*}
m _ {\mathbf {q}} ^ {2} & \leq \chi^ {yy} ( \mathbf {k} + \mathbf {q} ) \, F ( \mathbf {k} )\\
& \leq \chi^ {yy} \left[ h m _ {\mathbf {q}} + S ( S + 1 ) \bar {J} | \mathbf {k} | ^ {2} \right] .
\tag{6.40}
\end{align*}

反解并对布里渊区求和：

\begin{equation*}
\frac {m _ {\mathbf {q}} ^ {2}}{( 2 \pi ) ^ {d}} \int_ {0} ^ {\bar {k}} \frac {d k \, k ^ {d - 1}}{h m _ {\mathbf {q}} + S ( S + 1 ) \bar {J} | \mathbf {k} | ^ {2}} \leq \mathcal {N} ^ {- 1} \sum_ {\mathbf {k}} \chi^ {yy} ( \mathbf {k} ).
\tag{6.41}
\end{equation*}

因此，若

\begin{equation*}
\mathcal {N} ^ {- 1} \sum_ {\mathbf {k}} \chi^ {yy} ( \mathbf {k} ) = C < \infty ,
\tag{6.42}
\end{equation*}

且体系处于一维或二维，则 $m_{\mathbf{q}}$ 必须随 $h$ 趋于零才能满足 (6.41)。此时即使 $T = 0$ 也没有长程序。

可以证明：若激发谱有能隙

\begin{equation*}
E _ {m} - E _ {0} > \Delta \, , \quad \forall m \neq 0 ,
\tag{6.43}
\end{equation*}

且 $\Delta$ 与 $h$、$\mathcal{N}$ 无关，则 Heisenberg 模型的基态必为无序。能隙使我们能用关联函数估计磁化率 (6.39)：

\begin{equation*}
\chi^ {yy} ( \mathbf {k} ) \leq \frac {2}{\mathcal {N} \Delta} \sum_ {m} \left| \langle 0 | S _ {\mathbf {k}} ^ {y} | m \rangle \right| ^ {2} = 2 \frac {S ^ {yy} ( \mathbf {k} )}{\Delta} ,
\tag{6.44}
\end{equation*}

而 $\mathcal{N}^{-1}\sum_{\mathbf{k}}S^{yy}\leq S(S + 1)$，条件 (6.42) 即被满足。有趣的是，$\Delta$ 在不等式 (6.14) 中扮演着 $T$ 的角色。

例如，众所周知反铁磁整数自旋链的激发谱展现“Haldane 能隙”。$T = 0$ 的 Mermin--Wagner 型论证蕴含这些模型的基态没有真长程序。事实上，非线性 {$\sigma$} 模型分析预言其关联在大距离上指数衰减（见 15.2 节）。

逆命题（无能隙激发蕴含长程序）不成立。例如一维自旋 1/2 Heisenberg 反铁磁体激发无能隙，但 $T = 0$ 无长程序。

\begin{exercises}

\item 证明：由 (6.3) 有

\begin{equation*}
\lim _ {h \to 0 ^ {+}} m _ {\mathbf {q}} = - \lim _ {h \to 0 ^ {-}} m _ {\mathbf {q}} ,
\tag{6.45}
\end{equation*}

故长程序要求零频磁化率发散：

\begin{equation*}
\lim _ {h, \omega \rightarrow 0 ^ {+}} \chi^ {zz} ( \mathbf {q}, \omega ) = \infty .
\tag{6.46}
\end{equation*}

反之是否成立？（即磁化率发散是否蕴含长程序？）

\item 利用破缺对称的定义 (6.4) 与 Marshall 定理，证明：在两个子晶格等大小的有限二分格子上，Heisenberg 反铁磁体不可能有任何破缺的自旋对称。

\item 证明自发破缺对称 (6.4) 蕴含真长程序 (6.6)。提示：用 Cauchy--Schwarz 不等式证明在有限序场 $h_{q}$ 下

\begin{equation*}
\langle \mathbf {S} _ {\mathbf {q}} \cdot \mathbf {S} _ {- \mathbf {q}} \rangle \geq \mathcal {N} ^ {2} \left| \mathbf {m} _ {\mathbf {q}} \right| ^ {2}.
\tag{6.47}
\end{equation*}

再证明对转动对称的哈密顿量，(6.47) 的 $|\mathbf{h}_{\mathbf{q}}| \to 0$ 极限存在并导致 (6.6)。

\item 证明 (6.47) 蕴含 (6.7)。提示：用 Cauchy--Schwarz 不等式证明

\begin{equation*}
\langle \mathbf {S} _ {i} \cdot \mathbf {S} _ {j} \rangle \leq S ( S + 1 ) ,
\tag{6.48}
\end{equation*}

并用反证法证明 (6.7)。

\end{exercises}

\begin{chapbib}
\item Bogoliubov 不等式导出于
  N. N. Bogoliubov, Physik. Abhandl. Sowjetunion \textbf{6}, 1, 113, 229 (1962)。
\item 用 Bogoliubov 不等式确立长程序不存在的首例，归功于 Hohenberg 关于低维超流的论文：
  P. C. Hohenberg, Phys. Rev. \textbf{158}, 383 (1967)。
\item Mermin 与 Wagner 定理证明于其经典短文
  N. D. Mermin and H. Wagner, Phys. Rev. Lett. \textbf{17}, 1133 (1966)。
\item 相对论场论的类似定理由
  S. Coleman, Comm. Math. Phys. \textbf{31}, 259 (1973)
  证明。
\end{chapbib}
